% Section 1: Introduction (label prefix int). The statements of this section are proved in the body;
% each statement names its body version. Continued in part 2.
\section{Introduction}
\label{sec:int}

%-------------------------------------------------------------------------------
\subsection{The setting}
\label{sec:int-setting}

Let the torus $(\bbC^*)^5$ act on $\bbP^4$ with weights $\lambda_0,\dots,\lambda_4$ on $H^0(\bbP^4,\cO(1))$. The \emph{formal quintic} of Lho and Pandharipande \cite[\S0.2, (5)]{LP18}, \cite[\S0.2, (5)]{LP20} is the $\cO(5)$-twisted equivariant Gromov--Witten theory of $\bbP^4$ at the specialisation $\lambda_i=\zeta^i$, $\zeta=\re^{2\pi\ri/5}$, (see \cite[\S3.1; \S3.2, (9)]{GJR18}; the resulting hypergeometric series is studied in \cite{ZZ08}): its invariants are $N_{g,d}=\int_{[\overline{\mathcal M}_g(\bbP^4,d)]^{\mathrm{vir}}}e(R\pi_*f^*\cO(5))$, defined by virtual localisation \cite{MR1666787} before the weights are specialised, and $F_g=\sum_{d\geq0}N_{g,d}\Qn^d$ (\cref{def:set-fq}). By the genus-zero mirror theorem the $J$-function at $t=T(q)H$ is $I(q,z)/C_0(q)$, with $C_0=\sum_d\frac{(5d)!}{(d!)^5}q^d$ and the mirror map $T=I_1/C_0$ (\cref{thm:set-mirror}), and $\Qn=\re^{T(q)}$. With $u:=1-5^5q$ and $L:=u^{-1/5}=1/s$, the quantum product is cyclic (\cref{prop:set-cyclicLR}); $q=0$ is $u=1$, and $u=0$ is the \emph{conifold point}.

By Teleman's reconstruction \cite[Theorem~(2) in (1.15)]{MR2917177}\draftnote{numbers from arXiv v3; journal numbering not compared} (see \cite[\S4.1, Theorem~4.1]{GJR18}) $F_g$ is the Givental--Teleman sum over the stable graphs of genus $g$ of Givental's symplectic $R$-matrix $R(z)=1+O(z)$ acting on the topological field theory $\TFT$, in the convention of \cite[Defs.~2.2, 2.13]{PPZ15} (\cref{def:set-rows,def:set-graphsum,thm:set-GT}); at $q=0$ quantum Riemann--Roch, in the form \cite[(17)]{GJR18}, \cite[(19)--(20)]{GJR17} (which rest on \cite[\S2.2, Theorem~1]{Giv01} for the tangent weights and on \cite[Corollary~1]{MR2276766} for the twist), gives
\[
R(q=0)=\diag\bigl(h_{\bw_0},\dots,h_{\bw_4}\bigr),\qquad h_{\bw}(z)=\exp\Bigl(-\sum_{k\geq1}\bk_k\,z^{2k-1}\sum_{a\in\bw}a^{1-2k}\Bigr),\qquad \bk_k=\frac{B_{2k}}{2k(2k-1)},
\]
with $\bw_i=\{\zeta^i-\zeta^j\}_{j\neq i}\cup\{-5\zeta^i\}$, the input (QRR) (\cref{thm:set-qrr}); $h:=h_{\bw_0}$. In the conifold expansion defined below, where $L=s^{-1}$, the lowered rows $\Qr^{(i)}_a(z)=\eta(H^a,R(z)\tilde e_i)$ grow with $L$, and the natural variable is $\zh:=L^4z=z/s^4$.

\emph{Conifold frames.} Let $f_0,f_1,f_2$ be the solutions $f_k=u^k+O(u^3)$ at $u=0$ of the Picard--Fuchs operator $\PFfour$ of the quintic, and $\varpi_1=f_1+\frac7{10}f_2=u+O(u^2)$ the vanishing period (\cref{lem:set-frobbasis}). A \emph{conifold frame} is a period $\varpi_0=f_0+\alpha\varpi_1+\beta f_2$, $(\alpha,\beta)\in\bbQ^2$, and $\tc:=\varpi_1/\varpi_0=u+O(u^2)$ (\cref{def:set-frames}). The row coefficients lie in $\Nring\tensor\bbQ(\zeta)$, $\Nring$ a differential ring generated by the periods, $L^{\pm1}$ and a formal $\log q$ (\cref{def:set-confexp,cor:gau-rowsNN}); the \emph{conifold expansion} $\Cexp_{\alpha,\beta}\colon\Nring\to\bbQ\lser{s}$ expands at $u=0$ in the frame, and an expanded quantity is \emph{regular at the conifold point} if it lies in $\bbQ\fps{u}$.

\emph{The model and the natural form.} The model is the Landau--Ginzburg potential $\Wv(p)=p^4/4+v/p$ with primitive form $\dd p$, whose five critical points $p^5=v$ coalesce into the pole $p=0$ as $v\to0$ (\cref{def:set-model}). Its rows, the normalised stationary-phase series of $\int p^{\nu_b}\re^{-\Wv(p)/z}\dd p$, $\nu=(0,-1,3,2,1)$, form the model $R$-matrix $\Mo(t)$, $t=z/\varsigma^4$, $\varsigma=v^{1/5}$ (\cref{prop:set-modrows,def:set-M0}). Let $\tau(u):=1/{}_2F_1(1,\frac45;\frac95;u)=\bigl(\sum_{n\geq0}\frac{4u^n}{5n+4}\bigr)^{-1}$ and $v(u):=u\,\tau(u)^{-5/4}$ (\cref{def:set-tau}). At $v=v(u)$ the canonical coordinates of the model are those of the formal quintic measured from the conifold point, and $z/\varsigma^4=\tau(u)\zh=-\frac54\zeta^iz/(u_i-u_i^{\mathrm c})$ (\cref{lem:set-canon,prop:gau-equalprod}); so $M:=\Mo(\tau(u)\zh)$ is the $R$-matrix of the model at $v(u)$ in the variable of the formal quintic. Writing the expanded rows as $\Qr=R'M$, $R':=\Qr M^{-1}$, gives the \emph{natural form} of the graph sum: the outer graph sum $R'.(M.\TFT)$ has the vertices $M.\TFT$ and the edges and leaves of $R'$ (\cref{def:set-outer}).

%-------------------------------------------------------------------------------
\subsection{The conifold gap}
\label{sec:int-gap}

The conifold gap is a prediction of the topological B-model: near a conifold point the genus-$g$ potential, expanded in the flat coordinate of the vanishing period, has the universal leading pole of the $c=1$ string at the self-dual radius and no other pole \cite{Bershadsky:1993cx} (see \cite[\S1]{HKQ09}), \cite{Ghoshal:1995wm}, \cite[\S2.3]{HK07}, \cite[(1.2)]{HKQ09}. For the quintic threefold it was left open after \cite{GJR18,CGL21,CGL26} (see \cite[Remark~1.6]{GJR18}), was an input of high-genus computations \cite{HKQ09,LR22}, and is proved in \cite[Theorem~2]{CGYZ26} by master-space localisation. It holds for the resolved conifold \cite[\S1.1, footnote~1]{Br25}, for local $\bbP^2$ in all genera \cite{Br25} (after a check to genus $7$ in \cite{CI21}, see \cite[\S1.2]{Br25}), and for the topological recursion of toric mirror curves at a generic one-node degeneration \cite{FCCZ26}.

The potentials of the formal quintic satisfy holomorphic anomaly equations \cite[Theorems~1 and~2, \S0.4]{LP20}. By \cite[Remark~1.9]{CGYZ26}, the approach of \cite{GJR18} reduces the gap of the quintic to two key statements, the conifold gap for the formal quintic and the conifold regularity of its correlators with special insertions, and the authors of \cite{CGYZ26} expect their method to prove the first as well. We prove, conditionally on the inputs of \cref{sec:int-inputs}, a graph-local statement for the formal quintic (Theorem~A); summing over the graphs gives its conifold gap in every genus $g\geq2$ (Corollary~B), with the constant of \cite[Conjecture~1.5]{CGYZ26} (\cref{rmk:int-constant}).

The proof of \cite{CGYZ26} also factorises an $R$-matrix, that of a master-space theory \cite[Proposition~3.1]{CGYZ26}; here the $R$-matrix of the formal quintic itself is factorised, through a universal local model. \draftnote{The authors decide whether to cite their companion notes and how to relate the note to the sentence of \cite[\S1.1]{Br25} on $R$-matrix methods.}

The normalisation at $q=0$ is carried to the conifold point by an identity of formal stationary phase (Theorem~Gamma), through a change of chart as in Givental's treatment of projective spaces \cite[\S10]{MR1901075}. The analogous identity for the Landau--Ginzburg integral of Guo, Janda and Ruan is stated without proof in \cite[Remark~7.2]{GJR18}, and the value at the orbifold point is \cite[Lemma~7.7, (37)]{GJR18}; for Gamma factors in quantum cohomology see also \cite[Lemma~4.20]{MR2553377} and \cite{GGI16}. \draftnote{Re-check the latest versions of arXiv:2609.13921 and arXiv:2509.19298 before submission.}

% Section 1, part 2: the statements. Continued in part 3.

%-------------------------------------------------------------------------------
\subsection{Results}
\label{sec:int-results}

For a leg-free stable graph $\Gamma$, a matrix series $X$ and a cohomological field theory $\Omega$, $\mathrm{Cont}_\Gamma[E_X,q_X;\Omega]$ is the contraction of the vertex tensors of $\Omega$ along the edges $E_X$ of $X$, with the leaf point $q_X$ of $X$, as in \cite[Defs.~2.2, 2.13]{PPZ15} when $X$ is symplectic with $X(0)=1$, and with the edge of \cref{def:set-outer} in general, including the factor $|\Aut\Gamma|^{-1}$ (\cref{def:set-graphsum,def:set-outer,def:flh-outer}). The \emph{vertex graph} has one vertex, of genus $g$, and no edge; its term contains the leaf decorations.

\begin{thm}[Theorem~A]\genusstatus{AI-proposed, unchecked}
\label{thm:int-main}
Let $g\geq2$ and let $(\alpha,\beta)\in\bbQ^2$ define a conifold frame (\cref{def:set-frames}), with vanishing period $\varpi_1=u+O(u^2)$ and $\tc=\varpi_1/\varpi_0$. Let $M:=\Mo(\tau(u)\zh)$, exactly in $u$, and $R':=\Qr M^{-1}$, where $\Qr$ is the matrix of the lowered rows of \cref{def:set-rows} expanded in the frame (\cref{def:set-confexp}). For a leg-free stable graph $\Gamma$ of genus $g$ let
\[
T_\Gamma:=\mathrm{Cont}_\Gamma\bigl[E_{R'},q_{R'};M.\TFT\bigr]
\]
be its outer term (\cref{def:set-outer}), where $\TFT$ is the topological field theory of the formal quintic (\cref{def:set-frob}), expanded in the frame. Then:
\begin{enumerate}
\item $T_\Gamma\in\bbQ\fps{u}$ for every $\Gamma$ with at least one edge;
\item for the vertex graph, $T_\Gamma\in c_g\tc^{2-2g}+\bbQ\fps{u}$, where
\[
c_g=\frac{(-1)^{g-1}B_{2g}}{2g(2g-2)\,5^{g-1}}=\frac{|B_{2g}|}{2g(2g-2)\,5^{g-1}} .
\]
\end{enumerate}
\end{thm}

\Cref{thm:int-main} is proved as \cref{thm:mai-main}; for one $g$ it is Property N($g$) of \cref{def:set-N}. Here $c_2=1/1200$, $c_3=1/25200$ (\cref{rmk:set-cg}), and the hatted terms are $G_\Gamma=u^{2g-2}T_\Gamma$ (\cref{lem:flh-hat}). The statement is about each graph separately; we do not know how to obtain it from a statement about the singular part of the sum (\cref{rmk:mai-graphlocal}).

\begin{cor}[Corollary~B: the conifold gap of the formal quintic]\genusstatus{AI-proposed, unchecked}
\label{cor:int-gap}
Let $g\geq2$ and $(\alpha,\beta)\in\bbQ^2$. The conifold expansion $F_g^{\mathrm{cf}}\in\bbQ\lser{s}$ of the genus-$g$ potential of the formal quintic in the frame $(\alpha,\beta)$ (\cref{def:mai-Fcf}) satisfies
\[
F_g^{\mathrm{cf}}\in c_g\tc^{2-2g}+\bbQ\fps{u} .
\]
Since $\bbQ\fps{u}=\bbQ\fps{\tc}$, the series $F_g^{\mathrm{cf}}$ has no fractional power of $u$, no term $\tc^k$ with $3-2g\leq k\leq-1$, and leading term $c_g\tc^{2-2g}$.
\end{cor}

Here $F_g^{\mathrm{cf}}$ is the image under $\Cexp_{\alpha,\beta}$ of the genus-$g$ graph sum written in $\Nring$, whose evaluation is $F_g$ by (T) and (QRR); by the group property of the action \cite[Def.~2.13, Prop.~2.14]{PPZ15} it is the sum of the $T_\Gamma$ (\cref{cor:mai-gap}).

The proof compares two normalisations of the $R$-matrix. The cyclic systems of the formal quintic and of the model at $v(u)$ are related by a gauge regular at $u=0$, unique up to a scalar factor $\cn(z)$ (Theorem~U, \cref{thm:gau-U}); the \emph{canonical gauge} $\Pc$ is fixed at the conifold point by the gauge identity $\Pc\YC(v(u))=\YCQ$ between the vanishing solutions (\cref{thm:gau-G,def:gau-can}), the rows $\Qr$ at $q=0$ by (QRR). Both $\Qr$ and $\Pc\Qr^{\mathrm m}(v(u))$, $\Qr^{\mathrm m}$ the model rows, are formal WKB solutions of the same system with the same exponentials, so
\[
\Qr\,\Dn=\Pc\,\Qr^{\mathrm m}(v(u)),\qquad \Dn=\diag(\Dn_0,\dots,\Dn_4),\qquad \Dn_i(z)=\Dn_0(\zeta^{-i}z),\qquad \Dn_0\in\bbQ\fps{z}
\]
(\cref{thm:gau-R0,cor:gau-colours}; that $\Dn_0$ has rational coefficients is part of \cref{thm:con-glue}).

\begin{thm}[Theorem~C: the normalisation]\genusstatus{AI-proposed, unchecked}
\label{thm:int-D}
In every conifold frame, $\Dn_0=1$ in $\bbQ\fps{z}$. Hence $\Dn=1$, and the window $\Fg{g}$, that is $\Dn_0\equiv1\bmod z^{3g-2}$ (\cref{def:gau-Fg}), holds for every $g\geq2$. Equivalently, $\Qr=\Pc\,\Qr^{\mathrm m}(v(u))$: the actual gauge is the canonical gauge, and the $R$-matrix of the formal quintic, normalised at $q=0$ by (QRR), is in the canonical normalisation at the conifold point.
\end{thm}

\Cref{thm:int-D} is \cref{thm:gam-D} with \cref{cor:mai-actual}; its consequences for the natural form are \cref{cor:mai-natfact}. It is reduced to one value at $q=0$ of a saddle series:

\begin{thm}[Theorem~Gamma]\genusstatus{AI-proposed, unchecked}
\label{thm:int-gamma}
Let $c(z;L)=\sum_{n\geq0}c_n(L)z^n$, $c_0(L)=1$, be the full-Stirling saddle series of \cref{def:gam-c}: the normalised stationary-phase series, divided by its constant term ($\SPh$ of \cref{def:fsp-wick}), at $\sigma=L$ of
\[
\Gamma\Bigl(1+\frac{5\sigma}{z}\Bigr)\prod_{j=0}^{4}\Gamma\Bigl(\frac{\zeta^j-\sigma}{z}\Bigr)(-q)^{\sigma/z}\,\dd\sigma,\qquad q=q(L)=\frac{1-L^{-5}}{5^5},
\]
with every Gamma factor replaced by its formal Stirling form (\cref{def:fsp-stirling}). Then every $c_n$ lies in $\AL=\bbQ[L,L^{-1},(L^5-1)^{-1}]$, and:
\begin{enumerate}
\item[(a)] no $c_n$ has a pole at $L=1$, that is, at $q=0$, and $c(z;1)=h(z)$;
\item[(b)] no $c_n$ has a pole at $L=0$, that is, at $q=\infty$, and
\[
c(z;0)=\Ginf(z):=\exp\Bigl(\sum_{l\geq1}\frac{(-1)^{l+1}B_{5l+1}(\frac15)}{l(5l+1)}\,z^{5l}\Bigr),
\]
the series given by \cite[Lemma~7.7, (37)]{GJR18} at $m=5$, $k=1$, $\lambda=1$ in the notation there.
\end{enumerate}
\end{thm}

\Cref{thm:int-gamma} is formed by parts (i), (ii) and the first assertion of (iv) of \cref{thm:gam-main}, with \cref{thm:gam-ginf}. Part (a) is proved by a local Mellin inversion (\cref{thm:gam-lmi}) from the formal stationary-phase lemmas of \cref{sec:fsp} alone; by $c_n(L)=L^{4n}\mathrm{ph}^{\mathrm{SP}}_n(L^{-5})$ (\cref{prop:gam-hatflat}), Theorem~P (\cref{thm:con-Pb}) and the glue (\cref{thm:con-glue}) it is equivalent to $\Dn=1$ (\cref{cor:gam-equiv,cor:con-A2}). Part (b) says that $c(z;L)$ takes at $L=0$ the series displayed on the right of \cite[Lemma~7.7, (37)]{GJR18}, read as a series in $z$, which is obtained there for the Landau--Ginzburg series of that paper (we do not claim that the two series coincide); it identifies the free coefficients of the canonical colour-$0$ function (\cref{cor:gam-free}) and is not used for \cref{thm:int-main}.

\begin{rmk}
\label{rmk:int-constant}
The constant $c_g$ is that of \cite[Conjecture~1.5]{CGYZ26} for the quintic threefold, in the same normalisation $\varpi_1=u+O(u^2)$; the two theories have different potentials, so nothing about the quintic threefold follows (\cref{sec:sgn-lit}, \cref{rmk:mai-constant}). \draftnote{Numbers from arXiv:2609.13921v2, read in the HTML rendering; to be re-read in the PDF.}
\end{rmk}

% Section 1, part 3: inputs, what is not claimed, strategy, plan of the note.

%-------------------------------------------------------------------------------
\subsection{Inputs}
\label{sec:int-inputs}

The results are conditional on the following inputs. Statement numbers refer to the arXiv versions named in the bibliography.
\begin{description}
\item[(Q)] \emph{Quoted theorems.} Topological recursion, applied only to the model curve $x=\Wv(p)$, $y=p$ and its rescalings $x=\kappa\Wv$ (\cref{sec:mod}): \cite[Def.~4.2, (4-9), with the kernel (4-3); Theorems~4.2 and~4.6; Theorem~5.1, (5-14), with \S5.3, (5-27)--(5-30)]{Eynard:2007kz}; the theorem \cite[Theorem~4.1, with (4-1)--(4-3)]{DuninBarkowski:2012bw}, with its graph weights \cite[\S2.3, (2-22)--(2-26)]{DuninBarkowski:2012bw} (the dilaton leaves (2-23) and the factor $1/|\Aut|$ of (2-26)), as recalled in \cite[\S2.6, (2-30)--(2-32), (2-34), (2-35)]{Dunin-Barkowski:2015caa}, with matrices in the left-action convention of \cite[Remark~2.1]{Dunin-Barkowski:2015caa} (upper index the row), and with the action of \cite[\S2.3]{Dunin-Barkowski:2015caa} read with the factor $1/|\Aut\Gamma|$ and with the dilaton leaves of \cite[Def.~2.13]{PPZ15}, that is, as the action of \cite[Defs.~2.2, 2.13]{PPZ15}; it is applied to $x=4\Wv$, $y=p$, whose critical points $p^5=v$ have the pairwise distinct critical values $5p_c^4$. Its hypotheses, namely symplecticity, \cite[(2-32)]{Dunin-Barkowski:2015caa} and the edge condition \cite[(4-1)]{DuninBarkowski:2012bw}, are proved for this curve in \cref{prop:mod-Rtr}; the form \cite[(2-33)]{Dunin-Barkowski:2015caa} of the edge condition is misprinted (it has $+\sum_k$ in place of $\delta_{ij}-\sum_k$) and is not used, nor are \cite[Prop.~4.3]{Dunin-Barkowski:2015caa} and the sentence preceding \cite[Theorem~8.1]{Dunin-Barkowski:2015caa}; and \cite[Theorem~7.1: (7.6), (7.7) with its factor $1/\hbar^2$, (7.9), (7.10), with the kernel (1.1) and the assumptions of \S1.1]{ABDKS25}. The $R$-matrix action: \cite[\S0.2, (1); Def.~0.3, in particular (iii); (5); Remark~2.1 and (6); Def.~2.2, with the convention stated after it (p.~20); Def.~2.5, with the convention stated after it; Prop.~2.12, with the formula (8) in its proof; Def.~2.13; Prop.~2.14]{PPZ15}; Props.~2.12 and~2.14 are used as statements over $\bbQ$ and transferred to the other coefficient rings in \cref{def:set-graphsum}; (8) enters only through \cref{lem:set-pushins}.
\item[(QRR)] \emph{The $R$-matrix along the line: flatness and the value at the large-radius point} (\cref{thm:set-qrr}). The $R$-matrix $R$ of (T) has coefficients in $\End(\bbH)\otimes\bbQ(\zeta)\fps q$ and satisfies: (\emph{flatness}) for every colour $i$ the column $R\,\tilde e_i\,\re^{u_i/z}$ solves the quantum differential equation $z\partial_Tv=H*v$ along the line (\cref{def:set-rows}~(b)); (\emph{value}) $R(q=0)=\diag(h_{\bw_0},\dots,h_{\bw_4})$ in the normalised canonical basis at $q=0$. Flatness is a theorem about the $R$-matrix of the cohomological field theory, not a convention; the convention (Givental's) is only that the legs of the graph sum carry $R^{-1}(\psi)$. Flatness is used in \cref{prop:set-cyclicLR}~(iii) and in the last sentence of \cref{cor:set-cyclicPF5}~(b), on which \cref{lem:gau-colrows}, \cref{thm:gau-R0}~(iii), \cref{cor:gau-rowsNN}, \cref{thm:con-glue} and \cref{lem:mai-Psi} rest. The value is used for every colour, in the form $\Phi_i(z,1)=h_{\bw_i}(z)$ (\cref{eq:set-qrr}): in \cref{lem:gau-colrows} and, for colour $0$, in the glue \cref{thm:con-glue}. Sources for flatness: \cite[Lemma~5.2, (16)]{GJR18}, which is \cref{prop:set-cyclicLR}~(iii) at $\lambda=1$; \cite[\S4.2, (12)--(13)]{GJR17}; in general \cite[Section~1, Proposition (a)--(c)]{Giv01} and \cite[Props.~7.7, 8.2~(ii), 8.4, 8.7~(ii)]{MR2917177}. Sources for the value: \cite[(17), in the proof of Lemma~5.2]{GJR18}; \cite[(19)--(20), in the proof of Lemma~4.6, at $t=0$]{GJR17}, which state it for this theory; \cite[\S2.2, Theorem~1 and the Remark after it]{Giv01}; \cite[Theorem~1, Corollary~1]{MR2276766}, which give the tangent factor and the Euler-class twist separately.
\item[(M)] \emph{The genus-zero mirror theorem} \cite[Theorem~11.6]{MR1408320}\draftnote{numbers from arXiv v2; IMRN numbering not compared}, in the Calabi--Yau case $n=4$, $r=1$, $l_1=5$, $\lambda'=0$, with the set-up of \cite[\S9]{MR1408320} and \cite[Prop.~11.4]{MR1408320}: at $\lambda_i=\zeta^i$, $J(T(q)H,z)=I(q,z)/C_0(q)$ with $T=I_1/C_0$ (\cref{thm:set-mirror}).
\item[(G0)] \emph{The genus-zero axioms of the twisted theory}: the genus-zero twisted classes form a cohomological field theory with unit for $\eta$ \cite[\S4.1]{GJR18}, \cite[Def.~0.3]{PPZ15}, polynomial and $S_5$-symmetric in $\lambda$, with the string, divisor and splitting axioms, the genus-zero topological recursion relation and the quantum differential equation \cite[\S5.2, the display before (14), of which (14) is the form in the $q$-direction]{GJR18} (\cref{thm:set-G0}).
\item[(T)] \emph{Teleman's reconstruction, with the axioms of the twisted theory in all genera} (\cref{thm:set-GT}). (i) At $\lambda_i=\zeta^i$ the classes $\Omega_{g,n}$ of \cref{def:set-fq}, in all genera, form a cohomological field theory with flat unit $1$ for $\eta$ \cite[\S4.1]{GJR18}, \cite[Def.~0.3]{PPZ15}, and satisfy the divisor equation: for $2g-2+n>0$ the part of degree $d$ of $\pi_*\Omega_{g,n+1}(\gamma_1,\dots,\gamma_n,H)$ is $d$ times that of $\Omega_{g,n}(\gamma_1,\dots,\gamma_n)$, where $\pi$ forgets the last point; so the theory at $\Qn=\re^{T(q)}$ is the theory at the point $t=T(q)H$ with $\Qn=1$. (ii) This theory, which is semisimple (\cref{cor:set-cyclicPF5}~(a)), is the graph sum of \cref{def:set-graphsum}, $\Omega=R.\TFT$, for an $R$-matrix $R$ with coefficients in $\End(\bbH)\otimes\bbQ(\zeta)\fps q$, symplectic and with $R|_{z=0}=1$ \cite[Theorem~(2) in (1.15)]{MR2917177}, \cite[\S4.1, Theorem~4.1]{GJR18}; this is the $R$-matrix of \cref{def:set-rows}, to which (QRR) refers. In \cref{thm:int-main,thm:int-D} (T) identifies that $R$-matrix, whose rows define $T_\Gamma$ and $\Dn$; in \cref{cor:int-gap} it identifies the total with $F_g$.
\end{description}
In addition, the formula \cite[(1.7)]{AC96}, with $\kappa_0=2g-2+n$ \cite[p.~5]{AC96}, is used in part (i) of the dilaton lemma \cref{lem:set-dilaton}. The definitions use \cite[\S0.2, (5)]{LP18}, \cite[\S0.2, (5), (10)]{LP20} and \cite{MR1666787}; \cite{ZZ08} is cited for the series of \cref{lem:set-periods} only. The inputs are read with the following definitions and conventions of their sources: the normalisation of the $J$-function \cite[Section~6, (10)]{MR2276766} in (M) (\cref{thm:set-mirror}); the free energies \cite[Def.~4.3, (4-10)]{Eynard:2007kz} (\cref{conv:mod-kernel,cor:mod-trF}) and the set-up \cite[(5-1), the assumption before (5-4), (5-13)]{Eynard:2007kz} of \cite[Theorem~5.1]{Eynard:2007kz} (\cref{thm:mod-var}); the recursion \cite[(2-26)--(2-29)]{Dunin-Barkowski:2015caa}, to which the recollection of \cite[Theorem~4.1]{DuninBarkowski:2012bw} in \cite[\S2.6]{Dunin-Barkowski:2015caa} refers (\cref{conv:mod-kernel}); and the reading of the integrals \cite[(2-31), (2-32)]{Dunin-Barkowski:2015caa} as formal Laplace series (\cref{lem:mod-laplace,prop:mod-Rtr}): termwise Gaussian integration of the holomorphic part, as in \cite[(3-57)]{DuninBarkowski:2012bw}, and the value $\mathsf G_{-2}=-\vartheta$ for the polar term $w^{-2}\dd w$ (cf.\ \cite[(3-64)]{DuninBarkowski:2012bw}); and the edge condition \cite[(4-1)]{DuninBarkowski:2012bw} read in the coordinates $w_i$ of \cite[\S2.6]{Dunin-Barkowski:2015caa}, that is, as \eqref{eq:mod-edge} (\cref{prop:mod-Rtr}~(iii)).

\emph{Where the inputs enter.} (M) and (G0) enter only in \cref{sec:set} and its appendix \cref{sec:set-lrproof} (\cref{thm:set-mirror,prop:set-cyclicLR,cor:set-cyclicPF5}); (QRR) at the places named above: flatness in \cref{prop:set-cyclicLR}~(iii) and \cref{cor:set-cyclicPF5}~(b), the value in \cref{lem:gau-colrows} and \cref{thm:con-glue}; (T) in \cref{thm:set-GT} and \cref{cor:mai-gap}; topological recursion only in \cref{sec:mod}; the details are in \cref{sec:mai-thm}. \Cref{sec:gau} uses no quoted input beyond (M), (G0) and (QRR), which enter through the rows of \cref{sec:set} (flatness, \cref{def:set-rows}~(b), in \cref{thm:gau-R0}~(iii)--(iv)) and in \cref{lem:gau-colrows} (the value; hence in \cref{cor:gau-rowsNN}~(iii) and \cref{cor:gau-colours}); the transport lemmas of \cref{sec:tra} use \cite[Defs.~2.2, 2.13]{PPZ15} and the dilaton lemma, and Theorem~R adds no input; Theorem~P(a) (\cref{thm:con-Pa}) and part (a) of \cref{thm:int-gamma} use no quoted input.

\emph{Conditionality.} Theorem~A holds modulo the quoted theorems (Q) of \cite{Eynard:2007kz}, \cite{DuninBarkowski:2012bw} (as recalled in \cite{Dunin-Barkowski:2015caa}), \cite{ABDKS25} and \cite{PPZ15} with the statement numbers above, the flatness of the $R$-matrix and its value at the large-radius point (QRR), the genus-$0$ mirror theorem (M) and the genus-$0$ axioms (G0) of the twisted theory, Teleman's reconstruction with the axioms of the twisted theory in all genera (T), and the formula \cite[(1.7)]{AC96}; every other step is proved in the note.

In particular Theorem~U (over every field of characteristic $0$, \cref{thm:gau-U}) and every lemma of the chain are proved here; Stirling's series enters only through its formal proof as the stationary phase of Euler's integral (\cref{lem:fsp-stirling}), and the implicit function theorem only in its formal version (\cref{lem:fsp-ift}). \Cref{cor:int-gap} uses, beyond \cref{thm:int-main}, only (T), part of \cite{PPZ15} (listed after \cref{cor:mai-gap}) and proved statements. \Cref{thm:int-D} uses the $R$-matrix that (T) identifies only through $R|_{z=0}=1$, its coefficient ring $\bbQ(\zeta)\fps q$ and (QRR), that is, flatness (\cref{def:set-rows}~(b)) and the value at $q=0$, and the genus-zero data of \cref{sec:set}, which rest on (M) and (G0). \Cref{thm:int-gamma} uses no quoted input; the identity \cite[24.4.1]{DLMF} of Bernoulli polynomials used in part (b) is derived from their generating function in the proof of \cref{lem:fsp-dlmf5118}. No finite computation enters a proof; the computations that accompanied the work are summarised in \cref{sec:evi}.

All other citations are further sources, comparisons, conventions (\cref{sec:sgn}) or context, and enter no proof; among them \cite[Corollary~1.2]{Hock25} (\cref{rmk:mod-hock}), \cite[Theorem~4.4, (4-23); Corollary~4.1, (4-26)--(4-29); (A-1), (A-5), (A-25); Theorem~4.7, (4-30); (5-15)]{Eynard:2007kz}, \cite[(26)]{LP20}, \cite[Corollary~11.7]{MR1408320} (a misprinted form of (M)), \cite[Props.~6.1--6.2]{GR19}, \cite[\S2.4.3, (17); \S2.4.4, (19)--(20)]{LP20} and \cite{CGYZ26}.

%-------------------------------------------------------------------------------
\subsection{What is not claimed}
\label{sec:int-notclaimed}

\begin{enumerate}
\item Nothing is claimed about the quintic threefold or its graph sums; the formal quintic is a different theory. The conifold gap of the quintic threefold is \cite[Theorem~2]{CGYZ26}.
\item Nothing is claimed for frames other than those of \cref{def:set-frames} and the rescalings of their two periods (\cref{sec:sgn-lit}), for $g=1$, or away from $u=0$.
\item Nothing is claimed about the analytic continuation of the series $F_g(q)$ to a neighbourhood of $u=0$. The conifold expansion $\Cexp_{\alpha,\beta}$ is a formal change of frame on the ring $\Nring$ (\cref{def:set-confexp}), applied to the graph-sum representative $F_g^{\Nring}$ of $F_g$ (\cref{def:mai-Fcf}); this is the sense of ``conifold expansion'' and of ``conifold gap'' in \cref{cor:int-gap}.
\item The second key statement of \cite[Remark~1.9]{CGYZ26} is not addressed (the Model Lemma, \cref{def:red-ML}, concerns the model), and \cref{cor:int-gap} is not claimed to be literally the first, which is named there but not stated.
\item No novelty is claimed for the formal stationary-phase lemmas of \cref{sec:fsp}, proved for completeness, for the formula of $\Ginf$ \cite[Lemma~7.7, (37)]{GJR18}, or for the change of chart \cite[\S10]{MR1901075}.
\item Topological recursion is not applied to a mirror curve of the quintic or of the formal quintic, and \cite{FCCZ26} is not used.
\end{enumerate}

%-------------------------------------------------------------------------------
\subsection{Strategy}
\label{sec:int-strategy}

\emph{Reduction to the model (\cref{sec:gau,sec:tra,sec:red}).} In a conifold frame the rows satisfy a cyclic system with one simple pole at $u=0$ (\cref{prop:gau-quintic}), from \cref{prop:set-cyclicLR} and the cyclic form of $\PFfour$ (\cref{thm:gau-cyc}); the model pulled back by $v(u)$ has the same product of coefficients, which fixes $\tau$ (\cref{prop:gau-equalprod}). Theorem~U gives the regular gauge between them, and the gauge identity singles out the symplectic gauge $\Pc$ (\cref{thm:gau-G,prop:gau-symp}). Its edges lie in $\sKo\tensor\sKo$ for a slice $\sK$ defined by the vanishing solution of the model (\cref{prop:tra-can}), and Theorem~R (\cref{thm:red-main,cor:red-can}) reduces the regularity of the edge terms of the flat sum of $\Pc.\Omega^{\mathrm{mod}}_{v(u)}$, and the vertex law, to the Model Lemma on $\sK$ (\cref{def:red-ML}).

\emph{The model (\cref{sec:mod}).} The model cohomological field theory is topological recursion on its curve (\cref{thm:mod-tr}); insertions from $\sK$ become legs with pairings polynomial in $v$ (\cref{prop:mod-LD,thm:mod-V}), and the free energies are $c_hv^{2-2h}$ (\cref{thm:mod-gap}); this is the Model Lemma (\cref{thm:mod-MLK}).

\emph{The normalisation (\cref{sec:flh,sec:con,sec:gam}).} The comparison theorem (\cref{thm:flh-main}) identifies $T_\Gamma$ with the flat term, graph by graph, if $\Dn_0\equiv1\bmod z^{3g-2}$; Theorem~P (\cref{thm:con-Pa,thm:con-Pb}) and the glue (\cref{thm:con-glue}) turn $\Dn=1$ into part (a) of \cref{thm:int-gamma} (\cref{thm:gam-D}). \Cref{thm:mai-main} assembles the parts.

%-------------------------------------------------------------------------------
\subsection{Plan of the note}
\label{sec:int-plan}

\Cref{sec:set} defines the objects and states (M), (G0), (QRR) and (T); \cref{sec:gau} proves Theorem~U, the gauge identity and symplecticity; \cref{sec:tra} proves the transport lemmas and \cref{sec:red} Theorem~R; \cref{sec:mod} proves the Model Lemma; \cref{sec:flh} compares the flat and the outer sums; \cref{sec:con} proves Theorem~P and the glue, and \cref{sec:gam} \cref{thm:int-gamma,thm:int-D}. \Cref{sec:mai} proves \cref{thm:int-main,cor:int-gap}, and \cref{sec:cmp} records the natural factorisation and universality. The appendices prove the large-radius cyclic structure (\cref{sec:set-lrproof}), the cyclic form of $\PFfour$ (\cref{sec:gau-cycproof}), the $u$-adic convergence lemma (\cref{sec:tra-convproof}) and the lemmas of formal stationary phase (\cref{sec:fsp}), and record the sign conventions (\cref{sec:sgn}) and the finite computations (\cref{sec:evi}).


\subsection*{Acknowledgements}

This work was supported by the UKRI grant UKRI2394 and by the Leverhulme Trust grant LIP-2023-007.


